% FILE: 05_evidence_receipts_and_gates.tex
% Project: prompts — Downstream Manufacturing Prompt Corpus

\section{Evidence, Receipts, and Gates}
\label{sec:prompts-evidence}

\subsection{Tests}
\label{sec:prompts-evidence-tests}

The corpus specifies four families of test artifacts, none of which have been
executed as of 2026-05-31. These are authorized specifications, not achieved results.

\textbf{Absence-proof fixtures} (\texttt{execution-plans/absence-proof-fixtures.md}):
TypeScript fixtures using \texttt{@ts-expect-error} directives in
\texttt{tests/absence-proofs/}, validated by a custom runner that wraps
\texttt{tsc --noEmit}. Three fixture categories are specified: cross-boundary
capability leakage (verifying that \texttt{SandboxedCapability} cannot be passed as
\texttt{UnrestrictedCapability}), typestate violation (verifying that
\texttt{WasmModule<Unverified>} cannot be passed to \texttt{linkModule}), and memory
safety constraints (verifying that \texttt{Pointer<number, "unchecked">} cannot be
passed to \texttt{view.read32}). Rust-side fixtures use the \texttt{trybuild} crate.
Success criterion: 100\% of defined illegal states manifest as compilation errors.

\textbf{Structural law compile-fail fixtures} (GAP\_008 from
\texttt{DOWNSTREAM\_COMPAT\_GAP\_CLOSE.md}): All fixtures in
\texttt{tests/ui/compile\_fail/} must fail on structural error codes E0308, E0599,
E0277, or E0080 (not E0425 or E0432). Each fixture must have a matching
\texttt{.stderr} file. A fixture failing for the wrong error code is treated as a
defect equivalent to a compile-fail fixture that accidentally passes.

\textbf{PM4Py benchmark tests} (\texttt{downstream\_pm4py\_benchmark\_comparison.md}):
Seven algorithm comparisons across four dataset scales (10K, 100K, 1M, 50K
pathological events). Success criteria include: Inductive Miner $\geq 20\times$
faster than PM4Py; A* alignment $\geq 22\times$ faster; Petri net soundness
$\geq 30\times$ faster; DFG mining $\geq 25\times$ faster. Correctness tolerances:
fitness within $\pm0.001$, alignment cost within $\pm0.01$, precision within $\pm0.01$.
These targets are success criteria, not achieved results; no benchmark receipts exist
in this directory.

\textbf{Audit mesh tests} (\texttt{downstream\_audit\_mesh\_expansion.md}):
Automated capability verification with type boundary test requirements. Specifications
are present; execution artifacts are absent.

\subsection{Receipts}
\label{sec:prompts-evidence-receipts}

One receipt-bearing artifact exists within the \texttt{prompts/} directory:

\textbf{\texttt{checkpoint\_\_downstream\_prompts\_complete.md}}: The master
verification checklist. It confirms completeness of twelve downstream directive files
and specifies that all markdown links must use absolute \texttt{file:///} paths
without backticks. The checklist is a verification specification, not a cryptographic
receipt in the BLAKE3/Ed25519 sense. It does not carry a \texttt{validator\_signature}
field.

No \texttt{ProcessIntelligenceVerificationReceipt} JSON artifacts, no BLAKE3 hashes,
and no Ed25519 signatures exist as materialized files within the \texttt{prompts/}
directory. The receipt schemas are specifications of artifacts to be manufactured by
downstream engineering work.

\subsection{Checkpoints}
\label{sec:prompts-evidence-checkpoints}

The corpus references one upstream checkpoint as its authorization gate:
\texttt{PROCESS\_INTELLIGENCE\_ALIVE\_001}. This checkpoint is cited in
\texttt{downstream\_wasm4pm\_refactor.md} as the precondition for invoking the
refactor mandate: ``This mandate is authorized for downstream execution ONLY after
the \texttt{PROCESS\_INTELLIGENCE\_ALIVE\_001} seal is produced. Do not invoke before
that seal exists. Invoking before the seal means acting on unverified research
claims.''

The checkpoint is not materialized within the \texttt{prompts/} directory. Its
existence must be confirmed by examining \texttt{receipts/} and
\texttt{checkpoints/} in the parent research program. This is the primary open
question for the corpus's authorization status.

\subsection{ALIVE/PARTIAL/BLOCKED Status}
\label{sec:prompts-evidence-status}

The status of this corpus is \textbf{PARTIAL} per the evidence examined.

\textbf{What is complete}:
\begin{itemize}
  \item All twelve directive files are present and carry explicit authority sources.
  \item Every completion criterion is expressed as a verifiable artifact type
    (compile-pass fixture, compile-fail fixture with \texttt{.stderr}, JSON schema,
    or named Rust type).
  \item Directive verdicts of ``READY FOR ENGINEERING'' with ``DOCTORAL THESIS
    (99\% specification completeness)'' confidence are dated 2026-05-31.
  \item The gap-closure sequence (GAP\_007 $\to$ GAP\_001 $\to$ GAP\_002; parallel
    GAP\_003/004/008; parallel GAP\_005/006) is formally specified with dependency
    analysis.
  \item The Chatman Equation bridge blueprint is complete with BLAKE3/Ed25519
    cryptographic linkage specified to the field level.
\end{itemize}

\textbf{What is absent}:
\begin{itemize}
  \item The \texttt{PROCESS\_INTELLIGENCE\_ALIVE\_001} seal is not materialized in
    this directory. The authorization gate is documented but its issuance is
    unconfirmed from within the corpus.
  \item No benchmark receipts exist. The PM4Py speedup targets (20--30$\times$) are
    success criteria, not achieved measurements.
  \item No compile-fail \texttt{.stderr} fixture files exist within the corpus.
    These are specified as required by GAP\_008 but not present here.
  \item ZKP infrastructure (SNARK/STARK proofs referenced in
    \texttt{zkp-execution-boundaries.md}) has no implementation in either
    \texttt{wasm4pm} or \texttt{wasm4pm-compat} and no gap document tracking it.
\end{itemize}
